\documentclass[10pt,a4paper]{amsart}
\usepackage[margin=19mm]{geometry}
\usepackage{amsmath,amssymb,mathtools}
\usepackage[scr=rsfs,cal=euler]{mathalfa}
\usepackage{xfrac}
\usepackage[unicode,hidelinks,bookmarks=false]{hyperref}
\newcommand{\ie}{i.e.\ }
\newcommand{\cf}{cf.\ }
\newcommand{\resp}{respectively\ }
\newcommand{\Subsection}{\subsection}
\providecommand{\nobreakdash}{\nobreak\hbox{-}}
\setlength{\emergencystretch}{2em}
\multlinegap0pt
\usepackage{amssymb}
\usepackage{dsfont}
\usepackage{xcolor}
\newtheorem{prop}{Proposition}
\newcommand{\Om}{\Omega}
\newcommand{\R}{\mathbb{R}}
\def \bb{\hbox}
\title{Compactness results for Sign-Changing Solutions of critical nonlinear elliptic equations of low energy}
\author{Hussein Cheikh Ali, Bruno Premoselli}
\date{Source: arXiv:2412.00817v1 (CC BY 4.0)}
\begin{document}
\maketitle

\noindent\emph{Source and licence.} Compactness results for Sign-Changing Solutions of critical nonlinear elliptic equations of low energy. Version arXiv:2412.00817v1 (CC BY 4.0), distributed by arXiv under the Creative Commons Attribution 4.0 International licence, http://creativecommons.org/licenses/by/4.0/, as recorded in the arXiv licence metadata for this exact version and shown on the abstract page https://arxiv.org/abs/2412.00817v1. Full licence notice: \texttt{licenses/arxiv-2412.00817-CC-BY-4.0.txt}.\par\smallskip

\noindent\textit{Editorial scope.} Counted unit: the article's prop lemme:estu0 (source0.tex lines 817--822). The article's complete original proof of it is delivered with it (source0.tex lines 823--852, one \texttt{proof} environment of 30 lines). No mathematics is written, repaired or re-derived: the statement and the proof are the article's own lines, in the article's own notation, and no retained line is substituted or re-flowed.  Retained uncounted as the source-local context the proof needs: lines 386--388.  Nothing was omitted from any retained span, so the package carries no omission record.  No retained line contains a citation, so the wrapper carries no bibliography and no bibliography entry.  No retained line contains a figure environment, an image-inclusion command or drawing code, so no image is delivered and the package declares no asset dependency.  Editorial formatting only, in addition: an AMS \texttt{amsart} preamble that restates the article's own declarations and macros used by the retained lines, namely the environments \texttt{prop} on separate counters, together with the standard packages those lines need.  The retained environments print in this document as Proposition 1 in order of appearance, whereas the article prints the same items with the numbering of its own full document.  The complete native source of the article as distributed in the arXiv e-print for this exact version is bundled unchanged as \texttt{sources/169422/source0.tex}.
\begin{equation}\label{estGlai}
	G_{\alpha}(y,x)\leq \frac{C}{|y-x|^{n-2}} \min\left\lbrace 1,\frac{d(y,\partial\Om)d(x,\partial\Om)}{|y-x|^2}\right\rbrace \bb{ for all }\, x,y \in \Om,\,  x\neq y,
\end{equation}

\begin{prop}\label{lemme:estu0}
	Let $U \subset \Om$ be an open set. There exists a constant $C(U)$ such that $\lim_{|U| \to 0} C(U) = 0$ and such that, for all $y\in \Om$ and for all $\alpha \ge 1$,
	\begin{equation}\label{eq:estu0}
		\int_{U} G_{\alpha}(y,x)\, dx\leq C(U)\,  d(y,\partial\Om).
	\end{equation}
\end{prop}

\begin{proof}
We let $C(U) = \sup_{y \in \Om} \int_{U} |x-y|^{1-n} \, dx$. Since $\Om$ is bounded and $y \mapsto |y|^{1-n} \in L^{1}_{loc}(\R^n)$ we have $C(U) \to 0$ as $|U| \to 0$ by absolute continuity of the integral. Using \eqref{estGlai} yields
\begin{equation}\label{sumI1+I2}
	\int_{U} G_{\alpha}(y,x)\, dx= 
	O\left(  I_{1}(y)+I_{2}(y)\right) 
\end{equation}
where we have let, for $i=1,2$, 
\begin{equation*}
	I_i(y):= \int_{U_i}\frac{1}{|y-x|^{n-2}} \min\Big\{1,\frac{d(y,\partial\Om)d(x,\partial\Om)}{|y-x|^2}\Big\}\, dx,
\end{equation*}
and
\begin{eqnarray*}
	U_1:=U \cap \Big\{|y-x|<\frac{d(y,\partial\Om)}{2}\Big\} \bb{ and }U_2:=U \cap \Big\{|y-x|>\frac{d(y,\partial\Om)}{2}\Big\}.
\end{eqnarray*}
When $x \in U_1$ we have $|y-x|<\frac{d(y,\partial\Om)}{2}$ so that
\begin{equation*}
\begin{aligned}
I_1(y)\le\int_{U_1}\frac{1}{|y-x|^{n-2}}\, & \le \frac{d(y,\partial\Om)}{2}\int_{U}\frac{1}{|y-x|^{n-1}}\, dx \\
&  \le \frac{C(U)}{2} d(y, \partial \Om). 
\end{aligned}
\end{equation*}
When $x\in U_2$ we get that $d(x,\partial \Om) \leq 3|y-x|$.  We then get that 
\begin{equation*}
\begin{aligned}
I_2(y) \le d(y,\partial\Om)\int_{U_2} \frac{d(x,\partial\Om)}{|y-x|^{n}}
					& \le 3 d(y,\partial\Om) \int_{U}\frac{1}{|y-x|^{n-1}}\, dx \\
					& \le 3 C(U) d(y, \partial \Om). 
\end{aligned}
\end{equation*}
Combining these estimates proves Proposition \ref{lemme:estu0}.\end{proof}

\end{document}
